\section{RQ3: khả thi khi thời gian đi trôi dạt}

\label{muc:trongdiem}

\subsection{So bốn họ cổng trên cùng bộ sinh}

Phép kiểm tầng 2 hỏi một câu hẹp: khi thông tin thời gian đi đổi, cổng nào có thể loại
luôn phương án giữ tuyến? Bốn cổng có cùng bộ sinh chèn và các kiểm tra vật lý. Cổng
$C$ tính phần do quyết định; $V$ tính biến động nhìn thấy; $M^{+}$ giữ hạn chót một
phía từ lời hứa đầu; $M^{\pm}$ giữ cả giới hạn sớm và muộn từ lời hứa đầu. Các dạng
$M$ là đối chứng về hình thức ràng buộc, không phải bản tái lập toàn bộ thuật toán
của Schulz và Pfeiffer hay một nghiên cứu cửa sổ thời gian khác.

Trong đợt này, hệ số thời gian đi tại thời điểm $t$ là
$[f(t_0+t)/f(t_0)]^k$, với $f$ lấy từ chuỗi theo ngày. $k=0$ giữ thế giới tĩnh;
$k=1$ dùng biến thiên của chuỗi; $k=2$ khuếch đại biến thiên một cách tổng hợp.
Giá trị hiện tại được dùng để chiếu tuyến, không đọc hệ số tương lai. Hệ số áp chung
cho mọi cung, nên đây chưa phải mô hình tắc đường khác nhau giữa các khu phố.

Phải tách \emph{tập rỗng do cổng cam kết} khỏi \emph{tập rỗng do điều kiện vật lý}.
Ở loại thứ nhất, các lựa chọn còn lại bị cổng thời gian loại, kể cả giữ tuyến.
Ở loại thứ hai, khách đã nhận không còn thỏa điều kiện đón hoặc hoàn tất của bộ
kiểm tra vật lý. Không có định lý nào về sổ hành vi biến một kế hoạch vật lý đã
hỏng thành khả thi. Nhận định ``cổng $C$ không làm rỗng tập'' vì vậy không được
viết thành ``hệ thống $C$ luôn chạy được''.

Phép so trên 20 ô chính tách các nguyên nhân dừng. Trong thế giới
tĩnh, $C$-30 không dừng vì cam kết, còn $V$-30 có bốn ô dừng theo nguyên nhân này.
Với biến thiên $k=1$, $M^{+}$-30 dừng do cam kết ở 19/20 ô; $C$-30 là 0/20.
Tuy nhiên, $C$-30 có 11/20 ô dừng vì cửa sổ đón vật lý. Đây là bằng chứng về
khả năng bảo toàn phương án giữ tuyến của \emph{cổng}, cùng một giới hạn rõ của
hệ thực thi nghiêm ngặt đang dùng.

Nới hạn chót làm số ô rỗng do cam kết của $M^{+}$ giảm từ 19 ở trần 30 giây,
xuống 13 ở 60 giây rồi 3 ở 120 giây. Số ô dừng vật lý tương ứng là 1, 5 và 11.
Một cổng bớt chặt có thể cho thực thi đi xa hơn rồi mới gặp vi phạm vật lý;
không được đọc những số này như ba ước lượng độc lập về tỉ lệ phục vụ.
Trên toàn 56 ô $C$ của các panel và hệ số đã xét, không có lần dừng do cam kết,
nhưng có 14 lần dừng vật lý. Những lần dừng sớm cũng làm so sánh dịch vụ chỉ trên
các lần chạy sống sót trở nên lệch; khóa luận không thực hiện phép so đó.

\subsection{Phản ví dụ và dự đoán đã niêm phong}

Panel sáng/tối bổ sung dùng ngày 12--13/11, mỗi ngày bốn mẫu; có tám ô cho mỗi
khung. Nó được tách khỏi panel dùng xây dựng quy tắc dự đoán tầng 2, nhưng vẫn
là bằng chứng thăm dò trong một miền nhỏ. Buổi sáng, $M^{+}$-30 rỗng do cam kết
ở 4/8 ô, $V$-30 ở 8/8, còn $C$-30 ở 0/8. Buổi tối, khi thời gian đi giảm,
$M^{+}$ không rỗng do cam kết ở các trần 30, 60 và 120 giây. Ngược lại,
$M^{\pm}$-30 rỗng ở 7/8 ô do giới hạn sớm, và $V$-30 rỗng ở 8/8.
Hình~\ref{hinh:cong-troidat} so quan sát với dự đoán đã niêm phong ở hai khung.

\begin{figure}[htbp]
\centering
\includegraphics[width=\linewidth]{hinh/tier2_gates_final.pdf}
\caption{Số ô rỗng do cổng cam kết trong panel phát triển tầng 2, tám ô mỗi khung,
trần 30 giây. Các cột so quan sát với dự đoán niêm phong; hình không đếm những
lần dừng vì điều kiện vật lý. $M$ trên hình là khoảng hai phía $M^{\pm}$.}
\label{hinh:cong-troidat}
\end{figure}

Kết quả buổi tối giúp phân biệt hai dạng $M$. Xe nhanh lên không làm hỏng một
hạn chót chỉ cấm muộn; nó có thể làm hỏng một khoảng đối xứng cấm cả sớm. Một
trần tổng biến động nhìn thấy còn tính cả những dịch chuyển về phía sớm. Vì vậy
không có căn cứ gộp mọi cổng theo kết quả thành một nhóm luôn kém khả thi như nhau.

Quy tắc dự đoán trên chuỗi thời gian đi được cố định trước các lần chạy cần dự đoán.
Đối chiếu với kết quả cuối cho 98/102 dự đoán xác định đúng ở panel chính,
31/32 ở panel sáng và 40/40 ở panel tối. Mẫu số đếm các trường hợp mà quy tắc
đưa ra một dự đoán xác định, không phải số episode độc lập, số khách hay số định lý.
Phần dự đoán yếu về vi phạm vật lý được ghi riêng; không gộp vào ba tỉ số trên.

Những dự đoán sai có ý nghĩa phương pháp: quy tắc dùng một quỹ đạo tham chiếu,
trong khi một nhánh có thể nhận khách khác và đi theo quỹ đạo khác. Một điều kiện
đúng trên quỹ đạo tham chiếu chưa là điều kiện trên mọi phương án. Do đó độ đúng
cao kiểm một quy tắc chẩn đoán trong miền này, không chứng minh giả thiết (H)
của quan hệ bao hàm giữa mọi lịch.

\subsection{Phục hồi thay đổi câu hỏi so sánh}


Đợt tầng 2 dừng khi điều kiện nghiêm ngặt bị phá. Để nghiên cứu dịch vụ sau biến
động, đợt sau cho tiếp tục có ghi nhận vi phạm và so hai luật phục hồi ở
Mục~\ref{muc:phuchoi}. Luật (a) đóng băng quyết định khi cổng bế tắc. Luật (b)
cho một miễn trừ giới hạn đối với trạng thái đã vượt ngưỡng: phương án mới không
được làm đại lượng đang vi phạm xấu hơn phương án giữ tuyến. Các điều kiện khác
vẫn được kiểm; miễn trừ không xóa sổ hay đổi lại lời hứa đầu.

Trên tám ô có tắc nghẽn của phép kiểm phục hồi, gồm hai ô theo chuỗi thực và sáu
ô tổng hợp, $C$-30 phục vụ 697 khách. $V$-30 với luật (a) phục vụ 666; đổi sang
luật (b) thì đạt 697. $M^{+}$-30 với luật (b) cũng đạt 697. Vì vậy chênh dịch vụ
của một cổng còn phụ thuộc cách xử lý khi cổng không còn khả thi. Không được
quy toàn bộ chênh của cấu hình đóng băng cho riêng định nghĩa ngân sách.

Các phép kiểm hợp đồng có 28 trường hợp ghép cặp, trong đó sáu trường hợp kiểm
âm có cùng hành động giữa hai luật. Tổng 49 hành động đổi tuyến có miễn trừ của
luật (b) được kiểm theo trạng thái giữ tuyến. Đây là kiểm tra hành vi cài đặt
trong các kịch bản đã dựng, không phải bằng chứng rằng phục hồi sẽ duy trì mọi
lời hứa cứng. Luật ``không làm xấu thêm'' và ý tưởng phục hồi đã có trong văn
liệu điều phối; khóa luận không nhận chúng là ý tưởng mới.

Thế giới của đợt sau dùng mốc cập nhật đúng $900$ giây, cùng lịch hệ số giữa các
nhánh, và chỉ cung cấp hệ số hiện tại. Trong đợt tầng 2 cũ, cập nhật xảy ra ở
sự kiện mô phỏng kế tiếp: độ trễ so mốc có trung vị 22 giây và cực đại 210 giây.
Các nhánh cũ có thể nhận bản cập nhật ở thời điểm khác nhau. Vì vậy khóa luận
giữ các kết quả tầng 2 như kiểm cơ chế của phiên bản đó, không ghép chúng với
đợt mới thành một ước lượng dịch vụ duy nhất.

\subsection{Lời hứa chặt trên giờ cao điểm thực}


Phép kiểm W07 dùng tám ô 07:00--09:00 của ngày 12--13/11, tổng 864 yêu cầu,
chỉ với biến thiên từ chuỗi thực và luật phục hồi (b). Đây là panel phát triển
đã được đọc nhiều lần, không phải tập kiểm tra mới. Nhánh $U$ phục vụ 665 khách.
Ở trần 30 giây, $C$, $M^{+}$ và $V$ đều phục vụ 660; ở 60 giây, $C$ và
$M^{+}$ đều phục vụ 661. Tuy vậy, chúng có nghĩa vụ và gánh phục hồi khác nhau.

Ở 30 giây, tỉ lệ khách vượt lời hứa \emph{theo thước đo của chính luật}, lấy
trung bình không trọng số theo tám ô, là 0\% với $C$, 24,0\% với $M^{+}$ và
29,2\% với $V$. Với $M^{+}$-60, tỉ lệ là 7,3\%. Giá trị trước làm tròn là
24,04\% và 7,25\%; 7,25\% khác với tỉ lệ gộp theo khách $47/661=7{,}11\%$.
Các số này không phải kết quả khảo sát về thất vọng của hành khách.

Đối với $M^{+}$, lời hứa là giờ trả thực tế không quá muộn so với thông báo đầu.
Trôi dạt có thể phá lời hứa ấy ngay cả khi không nhận thêm khách. Cổng $C$
giữ sổ hành vi dưới trần nhưng vẫn có khách đến muộn theo thước hạn chót. Nếu
chấm $C$-30 bằng thước hạn chót 30 giây, tỉ lệ vượt cũng là 24,0\% trên panel
này. Do đó ``$C$ không phá lời hứa của mình'' nói về một hợp đồng hẹp hơn,
không nói giờ đến của khách tốt hơn một cách tuyệt đối.

Tỉ lệ quyết định có miễn trừ trung bình ở trần 30 giây là 58,6\% cho $M^{+}$
và 63,2\% cho $V$, còn $C$ là 0\%. Ở trần 60 giây, $M^{+}$ vẫn có 35,5\%
quyết định có miễn trừ. Cùng số khách phục vụ nhưng phải thường xuyên bỏ qua
một phần cổng là khác biệt vận hành cần công bố. Nó không làm số khách bị vượt
hạn chót biến mất, cũng không biến miễn trừ thành một bảo đảm cứng.

\subsection{Khác biệt ở trần trung gian có đáng kể về dịch vụ?}

W07X mở các mức 90--240 giây sau khi đã xem W07. Kế hoạch được niêm phong trước
lần chạy mới, nhưng vùng trần được chọn từ kết quả phát triển. Phép so vẫn chỉ
có tám ô của hai ngày; 104 lần chạy mới cộng 80 lần được dùng lại không tạo
thành 184 mẫu độc lập. Bảng~\ref{bang:w07x} giữ các mức làm đổi cách đọc kết quả.

\begin{table}[htbp]
\centering
\caption{Số yêu cầu hoàn tất ở các trần chọn lọc của W07X, trên 864 yêu cầu.
Cột cuối đếm các ô có cùng toàn bộ chuỗi hành động của $C$ và $M^{+}$.}
\label{bang:w07x}
\begin{tabular}{@{}rrrrr@{}}
\toprule
\textbf{Trần (s)} & \textbf{$C$} & \textbf{$M^{+}$} & \textbf{$V$} & \textbf{Cùng hành động}\\
\midrule
30 & 660 & 660 & 660 & 8/8\\
90 & 664 & 663 & 663 & 7/8\\
120 & 666 & 664 & 664 & 6/8\\
150 & 666 & 666 & 666 & 8/8\\
240 & 665 & 665 & 665 & 8/8\\
\bottomrule
\end{tabular}
\end{table}

Một đối chiếu tại quyết định phân kỳ đầu tiên cho thấy đúng cơ chế dự kiến.
Ở ngày 12/11, mẫu 1, quyết định thứ 294, khách cũ đã có 79,0 giây thay đổi khi
giữ tuyến. Cách ghép mới thêm 68,262 giây do quyết định. Tổng 147,262 giây vượt
hạn chót 90 hoặc 120 giây, nhưng phần riêng do quyết định còn dưới các trần đó.
$C$ nhận cách ghép, còn $M^{+}$ chặn. Đây là giải thích tại một trạng thái đã
ghi, không phải ước lượng tác động nhân quả tổng thể của trôi dạt.

Khác biệt làm thêm đúng một khách trong ô đó. Ở trần 120 giây có thêm một ô
khác phân kỳ, nên tổng chênh là hai khách. Không mức nào đạt tiêu chí đã đặt
cho một lợi ích dịch vụ: ít nhất ba khách và ít nhất ba ô. Từ 150 giây trở lên,
các nhánh có cùng số khách trên panel. Sáu trong tám dự đoán W07X đúng; hai
dự đoán về phạm vi phân kỳ và thứ tự gánh phục hồi sai. Khóa luận giữ các
trường hợp phân kỳ như bằng chứng cơ chế, không dùng chúng để tuyên bố $C$
phục vụ nhiều hơn các cổng còn lại.

RQ3 vì vậy cần hai lớp trả lời. Sổ hành vi giữ được nghĩa vụ của mình trong
các trường hợp đã xét, miễn là trạng thái giữ tuyến và ngân sách đáp ứng giả
thiết. Các cổng theo kết quả có thể cần phục hồi nhiều hơn, nhưng khi được
phục hồi chúng có thể phục vụ cùng số khách. Còn giờ khách thực sự thấy hoặc
thực sự đến vẫn phụ thuộc phần thay đổi ngoài sổ hành vi. Mục sau kiểm một
biên nhìn thấy có tính đến phần đó.
